\documentclass[12pt]{article}
\usepackage{amsmath,amssymb}
\usepackage[margin=1in]{geometry}

\begin{document}
    
\noindent
\textbf{Name: Sindayen, Lorenzo Danilo V.}

\vspace{0.2cm}

\noindent
\textbf{APM1220 - Applied Multivariate Data Analysis}

\vspace{0.7cm}

\hrule

\vspace{0.5cm}

\hrule

\vspace{0.7cm}

\noindent
\textbf{INSTRUCTIONS:}

\vspace{0.1cm}

\noindent
\begin{itemize}
    \item Answer all questions completely and show all necessary steps, intermediate calculations, and mathematical proofs clearly to receive full credit.

    \item Where matrix operations are required, ensure that dimensions conform before performing any multiplications, additions, or inversions.

    \item Present all final matrix entries and numerical results in exact, simplified form (e.g., fractions or radicals) unless decimals are specifically requested.

    \item Clearly label your final answers for each sub-problem (a, b, c, etc.).
\end{itemize}

\vspace{0.3cm}

\noindent
\textbf{Problem 2.2 (5 points)}

\vspace{0.5cm}

\noindent
Given the matrices

\[
A=
\begin{bmatrix}
-1 & 3\\
4 & 2
\end{bmatrix},
\qquad
B=
\begin{bmatrix}
4 & -3\\
1 & -2\\
-2 & 0
\end{bmatrix},
\qquad
\text{and}
\qquad
C=
\begin{bmatrix}
5\\
-4\\
2
\end{bmatrix}
\]

\noindent
perform the indicated multiplications:

\vspace{0.3cm}

\noindent
\textbf{(a) \(5A\)}

\vspace{0.5cm}

\noindent
\textbf{Sol.}

\vspace{0.3cm}


\[
5A
=
5
\begin{bmatrix}
-1 & 3\\
4 & 2
\end{bmatrix}
\]

\[
=
\begin{bmatrix}
5(-1) & 5(3)\\
5(4) & 5(2)
\end{bmatrix}
\]

\[
=
\begin{bmatrix}
-5 & 15\\
20 & 10
\end{bmatrix}
\]

\noindent
Hence,

\[
\boxed{
5A=
\begin{bmatrix}
-5 & 15\\
20 & 10
\end{bmatrix}
}
\]

\vspace{1cm}

\noindent
\textbf{(b) \(BA\)}

\vspace{0.5cm}

\noindent
\textbf{Sol.}

\vspace{0.3cm}

\[
BA
=
\begin{bmatrix}
4 & -3\\
1 & -2\\
-2 & 0
\end{bmatrix}
\begin{bmatrix}
-1 & 3\\
4 & 2
\end{bmatrix}
\]

\[
BA
=
\begin{bmatrix}
4(-1)+(-3)(4) & 4(3)+(-3)(2)\\
1(-1)+(-2)(4) & 1(3)+(-2)(2)\\
(-2)(-1)+0(4) & (-2)(3)+0(2)
\end{bmatrix}
\]

\[
=
\begin{bmatrix}
-4-12 & 12-6\\
-1-8 & 3-4\\
2+0 & -6+0
\end{bmatrix}
\]

\[
=
\begin{bmatrix}
-16 & 6\\
-9 & -1\\
2 & -6
\end{bmatrix}
\]

\noindent
Hence,

\[
\boxed{
BA=
\begin{bmatrix}
-16 & 6\\
-9 & -1\\
2 & -6
\end{bmatrix}
}
\]

\vspace{1cm}

\noindent
\textbf{(c) \(A'B'\)}

\vspace{0.5cm}

\noindent
\textbf{Sol.}

\vspace{0.3cm}

\noindent
The transpose of \(A\) is

\[
A'
=
\begin{bmatrix}
-1 & 4\\
3 & 2
\end{bmatrix}
\]

\noindent
and the transpose of \(B\) is

\[
B'
=
\begin{bmatrix}
4 & 1 & -2\\
-3 & -2 & 0
\end{bmatrix}.
\]

\vspace{0.5cm}

\noindent
So,

\vspace{0.5cm}

\[
A'B'
=
\begin{bmatrix}
-1 & 4\\
3 & 2
\end{bmatrix}
\begin{bmatrix}
4 & 1 & -2\\
-3 & -2 & 0
\end{bmatrix}
\]

\[
=
\begin{bmatrix}
(-1)(4)+4(-3) &
(-1)(1)+4(-2) &
(-1)(-2)+4(0)
\\
3(4)+2(-3) &
3(1)+2(-2) &
3(-2)+2(0)
\end{bmatrix}
\]

\[
=
\begin{bmatrix}
-4-12 & -1-8 & 2+0\\
12-6 & 3-4 & -6+0
\end{bmatrix}
\]

\[
=
\begin{bmatrix}
-16 & -9 & 2\\
6 & -1 & -6
\end{bmatrix}
\]

Hence,

\[
\boxed{
A'B'
=
\begin{bmatrix}
-16 & -9 & 2\\
6 & -1 & -6
\end{bmatrix}
}
\]

\vspace{1cm}

\noindent
\textbf{(d) \(C'B\)}

\vspace{0.5cm}

\noindent
\textbf{Sol.}

\vspace{0.3cm}

\noindent
The transpose of \(C\) is

\[
C'
=
\begin{bmatrix}
5 & -4 & 2
\end{bmatrix}.
\]

\noindent
So,

\[
C'B
=
\begin{bmatrix}
5 & -4 & 2
\end{bmatrix}
\begin{bmatrix}
4 & -3\\
1 & -2\\
-2 & 0
\end{bmatrix}
\]

\[
=
\begin{bmatrix}
5(4)+(-4)(1)+2(-2)
&
5(-3)+(-4)(-2)+2(0)
\end{bmatrix}
\]

\[
=
\begin{bmatrix}
20-4-4 & -15+8+0
\end{bmatrix}
\]

\[
=
\begin{bmatrix}
12 & -7
\end{bmatrix}
\]

Hence,

\[
\boxed{
C'B=
\begin{bmatrix}
12 & -7
\end{bmatrix}
}
\]

\vspace{1cm}

\noindent
\textbf{(e) Is \(AB\) defined?}

\vspace{0.5cm}

\noindent
\textbf{Sol.}

\vspace{0.3cm}

\noindent
The dimensions of the matrices are

\[
A_{2\times2}
\qquad\text{and}\qquad
B_{3\times2}.
\]

\noindent
For the multiplication \(AB\) to be defined, the number of columns of \(A\) must equal the number of rows of \(B\).

\vspace{0.3cm}

\noindent
However,

\[
2\neq3.
\]

\vspace{0.7cm}

\noindent
Therefore, the dimensions do not conform for matrix multiplication.

\noindent
Hence,

\[
\boxed{
AB\text{ is not defined.}
}
\]

\vspace{1cm}

\hrule

\vspace{0.7cm}

\noindent
\textbf{Problem 2.4 (6 points)}

\vspace{0.5cm}

\noindent
When \(A^{-1}\) and \(B^{-1}\) exist, prove each of the following:

\vspace{0.5cm}

\noindent
\textbf{(a) \((A')^{-1}=(A^{-1})'\)}

\vspace{0.5cm}

\noindent
\textbf{Sol.}

\vspace{0.3cm}

\noindent
Since \(A^{-1}\) exists,

\[
AA^{-1}=I
\]

\noindent
Transpose both sides

\[
(AA^{-1})'=I'
\]

\noindent
Using the property

\[
(AB)'=B'A'
\]

\noindent
we have

\[
(A^{-1})'A'=I'
\]

\noindent
Since the transpose of the identity matrix is itself,

\[
I'=I
\]

\noindent
Thus,

\[
(A^{-1})'A'=I
\]

\noindent
Since

\[
A^{-1}A=I
\]

\noindent
taking the transpose gives

\[
(A^{-1}A)'=I'
\]

\noindent
Therefore,

\[
A'(A^{-1})'=I
\]

\noindent
Hence, \((A^{-1})'\) is the inverse of \(A'\) and,

\[
\boxed{
(A')^{-1}=(A^{-1})'
}
\]

\vspace{1cm}

\noindent
\textbf{(b) \((AB)^{-1}=B^{-1}A^{-1}\)}

\vspace{0.5cm}

\noindent
\textbf{Sol.}

\vspace{0.3cm}

\noindent
Consider the product

\[
(B^{-1}A^{-1})AB
\]

\noindent
Using the associative property of matrix multiplication

\[
(B^{-1}A^{-1})AB
=
B^{-1}(A^{-1}A)B
\]

\noindent
Since

\[
A^{-1}A=I
\]

\noindent
we have

\[
B^{-1}(A^{-1}A)B
=
B^{-1}IB
\]

\noindent
Since multiplication by the identity matrix does not change a matrix

\[
B^{-1}IB
=
B^{-1}B
\]

\noindent
Since

\[
B^{-1}B=I
\]

\noindent
we obtain

\[
(B^{-1}A^{-1})AB
=
I
\]

\noindent
Now consider the product in the opposite order

\[
AB(B^{-1}A^{-1})
\]

\noindent
Using the associative property

\[
AB(B^{-1}A^{-1})
=
A(BB^{-1})A^{-1}
\]

\noindent
Since

\[
BB^{-1}=I
\]

\noindent
we have

\[
A(BB^{-1})A^{-1}
=
AIA^{-1}
\]

\noindent
Since multiplication by the identity matrix does not change a matrix

\[
AIA^{-1}
=
AA^{-1}
\]

\noindent
Since

\[
AA^{-1}=I
\]

\noindent
we obtain

\[
AB(B^{-1}A^{-1})
=
I
\]

\noindent
Therefore \(B^{-1}A^{-1}\) is the inverse of \(AB\)

\vspace{0.7cm}

\noindent
Hence

\[
\boxed{
(AB)^{-1}=B^{-1}A^{-1}
}
\]

\vspace{1cm}

\hrule

\vspace{0.7cm}

\noindent
\textbf{Problem 2.8 (6 points)}

\vspace{0.5cm}

\noindent
Given the matrix

\[
A=
\begin{bmatrix}
1 & 2\\
2 & -2
\end{bmatrix}
\]

\noindent
find the eigenvalues \(\lambda_1\) and \(\lambda_2\) and the associated normalized eigenvectors \(e_1\) and \(e_2\)

\noindent
Determine the spectral decomposition (2-16) of \(A\)

\vspace{0.5cm}

\noindent
\textbf{Sol.}

\vspace{0.3cm}

\noindent
To find the eigenvalues solve the characteristic equation

\[
|A-\lambda I|=0
\]

\[
\begin{vmatrix}
1-\lambda & 2\\
2 & -2-\lambda
\end{vmatrix}
=0
\]

\[
(1-\lambda)(-2-\lambda)-4=0
\]

\[
-2-\lambda+2\lambda+\lambda^2-4=0
\]

\[
\lambda^2+\lambda-6=0
\]

\[
(\lambda-2)(\lambda+3)=0
\]

\noindent
Hence,

\[
\boxed{
\lambda_1=2
\qquad
\lambda_2=-3
}
\]

\vspace{0.7cm}

\noindent
\textbf{Finding the eigenvector associated with \(\lambda_1=2\)}

\vspace{0.3cm}

\noindent
Solve

\[
(A-2I)e_1=0
\]

\[
A-2I
=
\begin{bmatrix}
1-2 & 2\\
2 & -2-2
\end{bmatrix}
\]

\[
=
\begin{bmatrix}
-1 & 2\\
2 & -4
\end{bmatrix}
\]

\noindent
Let

\[
e_1=
\begin{bmatrix}
x\\
y
\end{bmatrix}
\]

\noindent
Then

\[
\begin{bmatrix}
-1 & 2\\
2 & -4
\end{bmatrix}
\begin{bmatrix}
x\\
y
\end{bmatrix}
=
\begin{bmatrix}
0\\
0
\end{bmatrix}
\]

\noindent
Using the first equation

\[
-x+2y=0
\]

\[
x=2y
\]

\noindent
Let

\[
y=1
\]

\noindent
Then

\[
x=2
\]

\noindent
An eigenvector corresponding to \(\lambda_1=2\) is

\[
\begin{bmatrix}
2\\
1
\end{bmatrix}
\]

\noindent
The length of the eigenvector is

\[
\sqrt{2^2+1^2}
\]

\[
=\sqrt{5}
\]

\noindent
Hence,

\[
\boxed{
e_1=
\frac{1}{\sqrt{5}}
\begin{bmatrix}
2\\
1
\end{bmatrix}
}
\]

\vspace{0.7cm}

\noindent
\textbf{Finding the eigenvector associated with \(\lambda_2=-3\)}

\vspace{0.3cm}

\noindent
Solve

\[
(A+3I)e_2=0
\]

\[
A+3I
=
\begin{bmatrix}
1+3 & 2\\
2 & -2+3
\end{bmatrix}
\]

\[
=
\begin{bmatrix}
4 & 2\\
2 & 1
\end{bmatrix}
\]

\noindent
Let

\[
e_2=
\begin{bmatrix}
x\\
y
\end{bmatrix}
\]

\noindent
Then

\[
\begin{bmatrix}
4 & 2\\
2 & 1
\end{bmatrix}
\begin{bmatrix}
x\\
y
\end{bmatrix}
=
\begin{bmatrix}
0\\
0
\end{bmatrix}
\]

\noindent
Using the first equation

\[
4x+2y=0
\]

\[
2x+y=0
\]

\[
y=-2x
\]

\noindent
Let

\[
x=1
\]

\noindent
Then

\[
y=-2
\]

\noindent
An eigenvector corresponding to \(\lambda_2=-3\) is

\[
\begin{bmatrix}
1\\
-2
\end{bmatrix}
\]

\noindent
The length of the eigenvector is

\[
\sqrt{1^2+(-2)^2}
\]

\[
=\sqrt{5}
\]

\noindent
Hence,

\[
\boxed{
e_2=
\frac{1}{\sqrt{5}}
\begin{bmatrix}
1\\
-2
\end{bmatrix}
}
\]

\vspace{0.7cm}

\noindent
\textbf{Finding the spectral decomposition}

\vspace{0.3cm}

\noindent
Since \(A\) is symmetric its spectral decomposition is

\[
A
=
\lambda_1e_1e_1'
+
\lambda_2e_2e_2'
\]

\noindent
Substituting the eigenvalues and normalized eigenvectors

\[
A
=
2e_1e_1'
-
3e_2e_2'
\]

\noindent
First calculate \(e_1e_1'\)

\[
e_1e_1'
=
\frac{1}{\sqrt{5}}
\begin{bmatrix}
2\\
1
\end{bmatrix}
\frac{1}{\sqrt{5}}
\begin{bmatrix}
2 & 1
\end{bmatrix}
\]

\[
=
\frac{1}{5}
\begin{bmatrix}
2\\
1
\end{bmatrix}
\begin{bmatrix}
2 & 1
\end{bmatrix}
\]

\[
=
\frac{1}{5}
\begin{bmatrix}
4 & 2\\
2 & 1
\end{bmatrix}
\]

\noindent
Now calculate \(e_2e_2'\)

\[
e_2e_2'
=
\frac{1}{\sqrt{5}}
\begin{bmatrix}
1\\
-2
\end{bmatrix}
\frac{1}{\sqrt{5}}
\begin{bmatrix}
1 & -2
\end{bmatrix}
\]

\[
=
\frac{1}{5}
\begin{bmatrix}
1\\
-2
\end{bmatrix}
\begin{bmatrix}
1 & -2
\end{bmatrix}
\]

\[
=
\frac{1}{5}
\begin{bmatrix}
1 & -2\\
-2 & 4
\end{bmatrix}
\]

\noindent
Therefore,

\[
A
=
2
\left(
\frac{1}{5}
\begin{bmatrix}
4 & 2\\
2 & 1
\end{bmatrix}
\right)
-
3
\left(
\frac{1}{5}
\begin{bmatrix}
1 & -2\\
-2 & 4
\end{bmatrix}
\right)
\]

\[
=
\frac{1}{5}
\begin{bmatrix}
8 & 4\\
4 & 2
\end{bmatrix}
-
\frac{1}{5}
\begin{bmatrix}
3 & -6\\
-6 & 12
\end{bmatrix}
\]

\[
=
\frac{1}{5}
\left[
\begin{bmatrix}
8 & 4\\
4 & 2
\end{bmatrix}
-
\begin{bmatrix}
3 & -6\\
-6 & 12
\end{bmatrix}
\right]
\]

\[
=
\frac{1}{5}
\begin{bmatrix}
5 & 10\\
10 & -10
\end{bmatrix}
\]

\[
=
\begin{bmatrix}
1 & 2\\
2 & -2
\end{bmatrix}
\]

\noindent
Hence,

\[
\boxed{
A
=
2e_1e_1'
-
3e_2e_2'
}
\]

\noindent
where

\[
\boxed{
e_1=
\frac{1}{\sqrt{5}}
\begin{bmatrix}
2\\
1
\end{bmatrix}
\qquad
e_2=
\frac{1}{\sqrt{5}}
\begin{bmatrix}
1\\
-2
\end{bmatrix}
}
\]

\vspace{1cm}

\hrule

\vspace{0.7cm}

\noindent
\textbf{Problem 2.11 (5 points)}

\vspace{0.5cm}

\noindent
Show that the determinant of the \(p\times p\) diagonal matrix \(A=\{a_{ij}\}\) with \(a_{ij}=0\) for \(i\neq j\) is given by the product of the diagonal elements

\[
|A|=a_{11}a_{22}\cdots a_{pp}
\]

\vspace{0.5cm}

\noindent
\textbf{Sol.}

\vspace{0.3cm}

\noindent
Consider the \(p\times p\) diagonal matrix

\[
A=
\begin{bmatrix}
a_{11} & 0 & 0 & \cdots & 0\\
0 & a_{22} & 0 & \cdots & 0\\
0 & 0 & a_{33} & \cdots & 0\\
\vdots & \vdots & \vdots & \ddots & \vdots\\
0 & 0 & 0 & \cdots & a_{pp}
\end{bmatrix}
\]

\noindent
By the definition of the determinant expand along the first row

\[
|A|
=
a_{11}A_{11}
+
0
+\cdots+
0
\]

\[
|A|
=
a_{11}A_{11}
\]

\noindent
Delete the first row and first column

\[
A_{11}
=
\begin{vmatrix}
a_{22} & 0 & \cdots & 0\\
0 & a_{33} & \cdots & 0\\
\vdots & \vdots & \ddots & \vdots\\
0 & 0 & \cdots & a_{pp}
\end{vmatrix}
\]

\noindent
Expand along the first row

\[
A_{11}
=
a_{22}A_{22}
\]

\noindent
Substitute into the determinant of \(A\)

\[
|A|
=
a_{11}a_{22}A_{22}
\]

\[
|A|
=
a_{11}a_{22}a_{33}\cdots a_{pp}
\]

\noindent
Hence, the determinant of a diagonal matrix is the product of its diagonal elements

\[
\boxed{
|A|
=
a_{11}a_{22}\cdots a_{pp}
}
\]

\vspace{1cm}

\hrule

\vspace{0.7cm}

\noindent
\textbf{Problem 2.24 (5 points)}

\vspace{0.5cm}

\noindent
Let \(X\) have covariance matrix

\[
\Sigma
=
\begin{bmatrix}
4 & 0 & 0\\
0 & 9 & 0\\
0 & 0 & 1
\end{bmatrix}
\]

\noindent
Find:

\begin{flalign*}
\textbf{(a) } \Sigma^{-1} &&
\end{flalign*}

\begin{flalign*}
\textbf{(b) The eigenvalues and eigenvectors of } \Sigma &&
\end{flalign*}

\begin{flalign*}
\textbf{(c) The eigenvalues and eigenvectors of } \Sigma^{-1} &&
\end{flalign*}

\vspace{1.0cm}

\noindent
\textbf{Answers: }

\vspace{1.0cm}

\noindent
\textbf{(a) \(\Sigma^{-1}\)}

\vspace{0.5cm}

\noindent
\textbf{Sol.}

\vspace{0.3cm}

\noindent
Since \(\Sigma\) is a diagonal matrix its inverse is obtained by taking the reciprocal of each diagonal element

\vspace{0.3cm}

\noindent
Hence,

\[
\boxed{
\Sigma^{-1}
=
\begin{bmatrix}
\frac{1}{4} & 0 & 0\\
0 & \frac{1}{9} & 0\\
0 & 0 & 1
\end{bmatrix}
}
\]

\vspace{1cm}

\noindent
\textbf{(b) The eigenvalues and eigenvectors of \(\Sigma\)}

\vspace{0.5cm}

\noindent
\textbf{Sol.}

\vspace{0.3cm}

\noindent
Solve for

\[
|\Sigma-\lambda I|=0
\]

\[
\begin{vmatrix}
4-\lambda & 0 & 0\\
0 & 9-\lambda & 0\\
0 & 0 & 1-\lambda
\end{vmatrix}
=
0
\]

\[
(4-\lambda)(9-\lambda)(1-\lambda)=0
\]

\noindent
Hence, the eigenvalues are

\[
\boxed{
\lambda_1=4
\qquad
\lambda_2=9
\qquad
\lambda_3=1
}
\]

\vspace{0.5cm}

\noindent
\textbf{Eigenvector corresponding to \(\lambda_1=4\)}

\vspace{0.3cm}

\noindent
Solve for

\[
(\Sigma-4I)e_1=0
\]

\[
\begin{bmatrix}
0 & 0 & 0\\
0 & 5 & 0\\
0 & 0 & -3
\end{bmatrix}
\begin{bmatrix}
x\\
y\\
z
\end{bmatrix}
=
\begin{bmatrix}
0\\
0\\
0
\end{bmatrix}
\]

\noindent
This gives

\[
5y=0
\]

\[
-3z=0
\]

\[
y=0
\qquad
z=0
\]

\noindent
Let

\[
x=1
\]

\noindent
Hence,

\[
\boxed{
e_1=
\begin{bmatrix}
1\\
0\\
0
\end{bmatrix}
}
\]

\vspace{0.5cm}

\noindent
\textbf{Eigenvector corresponding to \(\lambda_2=9\)}

\vspace{0.3cm}

\noindent
Solve for

\[
(\Sigma-9I)e_2=0
\]

\[
\begin{bmatrix}
-5 & 0 & 0\\
0 & 0 & 0\\
0 & 0 & -8
\end{bmatrix}
\begin{bmatrix}
x\\
y\\
z
\end{bmatrix}
=
\begin{bmatrix}
0\\
0\\
0
\end{bmatrix}
\]

\noindent
This gives

\[
-5x=0
\]

\[
-8z=0
\]

\[
x=0
\qquad
z=0
\]

\noindent
Let

\[
y=1
\]

\noindent
Hence,

\[
\boxed{
e_2=
\begin{bmatrix}
0\\
1\\
0
\end{bmatrix}
}
\]

\vspace{0.5cm}

\noindent
\textbf{Eigenvector corresponding to \(\lambda_3=1\)}

\vspace{0.3cm}

\noindent
Solve for

\[
(\Sigma-I)e_3=0
\]

\[
\begin{bmatrix}
3 & 0 & 0\\
0 & 8 & 0\\
0 & 0 & 0
\end{bmatrix}
\begin{bmatrix}
x\\
y\\
z
\end{bmatrix}
=
\begin{bmatrix}
0\\
0\\
0
\end{bmatrix}
\]

\noindent
This gives

\[
3x=0
\]

\[
8y=0
\]

\[
x=0
\qquad
y=0
\]

\noindent
Let

\[
z=1
\]

\noindent
Hence,

\[
\boxed{
e_3=
\begin{bmatrix}
0\\
0\\
1
\end{bmatrix}
}
\]

\vspace{0.5cm}

\noindent
Hence, the eigenvalues and corresponding eigenvectors of \(\Sigma\) are

\[
\boxed{
\lambda_1=4
\qquad
e_1=
\begin{bmatrix}
1\\
0\\
0
\end{bmatrix}
}
\]

\[
\boxed{
\lambda_2=9
\qquad
e_2=
\begin{bmatrix}
0\\
1\\
0
\end{bmatrix}
}
\]

\[
\boxed{
\lambda_3=1
\qquad
e_3=
\begin{bmatrix}
0\\
0\\
1
\end{bmatrix}
}
\]

\vspace{1cm}

\noindent
\textbf{(c) The eigenvalues and eigenvectors of \(\Sigma^{-1}\)}

\vspace{0.5cm}

\noindent
\textbf{Sol.}

\vspace{0.3cm}

\[
\Sigma^{-1}
=
\begin{bmatrix}
\frac{1}{4} & 0 & 0\\
0 & \frac{1}{9} & 0\\
0 & 0 & 1
\end{bmatrix}
\]

\noindent
Since the eigenvalues of the inverse of a matrix are the reciprocals of the nonzero eigenvalues of the original matrix

\[
\lambda_1^*
=
\frac{1}{4}
\]

\[
\lambda_2^*
=
\frac{1}{9}
\]

\[
\lambda_3^*
=
1
\]

\noindent
The eigenvectors remain unchanged because

\[
\Sigma e_i=\lambda_i e_i
\]

\[
\Sigma^{-1}e_i
=
\frac{1}{\lambda_i}e_i
\]

\noindent
Hence,

\[
\boxed{
\lambda_1^*=\frac{1}{4}
\qquad
e_1=
\begin{bmatrix}
1\\
0\\
0
\end{bmatrix}
}
\]

\[
\boxed{
\lambda_2^*=\frac{1}{9}
\qquad
e_2=
\begin{bmatrix}
0\\
1\\
0
\end{bmatrix}
}
\]

\[
\boxed{
\lambda_3^*=1
\qquad
e_3=
\begin{bmatrix}
0\\
0\\
1
\end{bmatrix}
}
\]

\vspace{1cm}

\hrule

\vspace{0.7cm}

\noindent
\textbf{Problem 2.30 (8 points)}

\vspace{0.5cm}

\noindent
You are given the random vector

\[
X'
=
[X_1,X_2,X_3,X_4]
\]

\noindent
with mean vector

\[
\mu_X'
=
[4,3,2,1]
\]

\noindent
and variance-covariance matrix

\[
\Sigma_X
=
\begin{bmatrix}
3 & 0 & 2 & 2\\
0 & 1 & 1 & 0\\
2 & 1 & 9 & -2\\
2 & 0 & -2 & 4
\end{bmatrix}
\]

\noindent
Partition \(X\) as

\[
X
=
\begin{bmatrix}
X_1\\
X_2\\
X_3\\
X_4
\end{bmatrix}
=
\begin{bmatrix}
X^{(1)}\\
X^{(2)}
\end{bmatrix}
\]

\noindent
Let

\[
A=
\begin{bmatrix}
1 & 2
\end{bmatrix}
\]

\noindent
and

\[
B=
\begin{bmatrix}
1 & -2\\
2 & -1
\end{bmatrix}
\]

\noindent
and consider the linear combinations \(AX^{(1)}\) and \(BX^{(2)}\)

\vspace{0.5cm}

\noindent
\textbf{(a) \(E(X^{(1)})\)}

\vspace{0.5cm}

\noindent
\textbf{Sol.}

\vspace{0.3cm}

\noindent
The first two elements of the mean vector form the mean vector of \(X^{(1)}\)

\vspace{0.5cm}

\noindent
Hence,

\[
\boxed{
E(X^{(1)})
=
\begin{bmatrix}
4\\
3
\end{bmatrix}
}
\]

\vspace{1cm}

\noindent
\textbf{(b) \(E(AX^{(1)})\)}

\vspace{0.5cm}

\noindent
\textbf{Sol.}

\vspace{0.3cm}

\noindent
Using the property

\[
E(AX^{(1)})
=
AE(X^{(1)})
\]

\[
E(AX^{(1)})
=
\begin{bmatrix}
1 & 2
\end{bmatrix}
\begin{bmatrix}
4\\
3
\end{bmatrix}
\]

\[
=
1(4)+2(3)
\]

\[
=
4+6
\]

\[
=
10
\]

\noindent
Hence,

\[
\boxed{
E(AX^{(1)})=10
}
\]

\vspace{1cm}

\noindent
\textbf{(c) \(\operatorname{Cov}(X^{(1)})\)}

\vspace{0.5cm}

\noindent
\textbf{Sol.}

\vspace{0.3cm}

\noindent
The covariance matrix of \(X^{(1)}\) is the upper-left \(2\times2\) submatrix of \(\Sigma_X\)

\[
\boxed{
\operatorname{Cov}(X^{(1)})
=
\begin{bmatrix}
3 & 0\\
0 & 1
\end{bmatrix}
}
\]

\vspace{1cm}

\noindent
\textbf{(d) \(\operatorname{Cov}(AX^{(1)})\)}

\vspace{0.5cm}

\noindent
\textbf{Sol.}

\vspace{0.3cm}

\noindent
Using the property

\[
\operatorname{Cov}(AX^{(1)})
=
A\operatorname{Cov}(X^{(1)})A'
\]

\[
\operatorname{Cov}(AX^{(1)})
=
\begin{bmatrix}
1 & 2
\end{bmatrix}
\begin{bmatrix}
3 & 0\\
0 & 1
\end{bmatrix}
\begin{bmatrix}
1\\
2
\end{bmatrix}
\]

\[
=
\begin{bmatrix}
1 & 2
\end{bmatrix}
\begin{bmatrix}
3 & 0\\
0 & 2
\end{bmatrix}
\]

\[
=
\begin{bmatrix}
3 & 4
\end{bmatrix}
\begin{bmatrix}
1\\
2
\end{bmatrix}
\]

\[
=
3(1)+4(2)
\]

\[
=
3+8
\]

\[
=
11
\]

\noindent
Hence,

\[
\boxed{
\operatorname{Cov}(AX^{(1)})=11
}
\]

\vspace{1cm}

\noindent
\textbf{(e) \(E(X^{(2)})\)}

\vspace{0.5cm}

\noindent
\textbf{Sol.}

\vspace{0.3cm}

\noindent
The last two elements of the mean vector form the mean vector of \(X^{(2)}\)

\vspace{0.3cm}

\noindent
Hence,

\vspace{0.3cm}

\[
\boxed{
E(X^{(2)})
=
\begin{bmatrix}
2\\
1
\end{bmatrix}
}
\]

\vspace{1cm}

\noindent
\textbf{(f) \(E(BX^{(2)})\)}

\vspace{0.5cm}

\noindent
\textbf{Sol.}

\vspace{0.3cm}

\noindent
Using the property

\[
E(BX^{(2)})
=
BE(X^{(2)})
\]

\[
E(BX^{(2)})
=
\begin{bmatrix}
1 & -2\\
2 & -1
\end{bmatrix}
\begin{bmatrix}
2\\
1
\end{bmatrix}
\]

\[
=
\begin{bmatrix}
1(2)+(-2)(1)\\
2(2)+(-1)(1)
\end{bmatrix}
\]

\[
=
\begin{bmatrix}
2-2\\
4-1
\end{bmatrix}
\]

\[
=
\begin{bmatrix}
0\\
3
\end{bmatrix}
\]

\noindent
Hence,

\[
\boxed{
E(BX^{(2)})
=
\begin{bmatrix}
0\\
3
\end{bmatrix}
}
\]

\vspace{1cm}

\noindent
\textbf{(g) \(\operatorname{Cov}(X^{(2)})\)}

\vspace{0.5cm}

\noindent
\textbf{Sol.}

\vspace{0.3cm}

\noindent
The covariance matrix of \(X^{(2)}\) is the lower-right \(2\times2\) submatrix of \(\Sigma_X\)

\vspace{0.3cm}

\noindent
Hence,

\vspace{0.3cm}

\[
\boxed{
\operatorname{Cov}(X^{(2)})
=
\begin{bmatrix}
9 & -2\\
-2 & 4
\end{bmatrix}
}
\]

\vspace{1cm}

\noindent
\textbf{(h) \(\operatorname{Cov}(BX^{(2)})\)}

\vspace{0.5cm}

\noindent
\textbf{Sol.}

\vspace{0.3cm}

\noindent
Using the property

\[
\operatorname{Cov}(BX^{(2)})
=
B\operatorname{Cov}(X^{(2)})B'
\]

\[
\operatorname{Cov}(BX^{(2)})
=
\begin{bmatrix}
1 & -2\\
2 & -1
\end{bmatrix}
\begin{bmatrix}
9 & -2\\
-2 & 4
\end{bmatrix}
\begin{bmatrix}
1 & 2\\
-2 & -1
\end{bmatrix}
\]

\[
B\operatorname{Cov}(X^{(2)})
=
\begin{bmatrix}
1(9)+(-2)(-2) & 1(-2)+(-2)(4)\\
2(9)+(-1)(-2) & 2(-2)+(-1)(4)
\end{bmatrix}
\]

\[
=
\begin{bmatrix}
9+4 & -2-8\\
18+2 & -4-4
\end{bmatrix}
\]

\[
=
\begin{bmatrix}
13 & -10\\
20 & -8
\end{bmatrix}
\]

\noindent
Therefore,

\[
\operatorname{Cov}(BX^{(2)})
=
\begin{bmatrix}
13 & -10\\
20 & -8
\end{bmatrix}
\begin{bmatrix}
1 & 2\\
-2 & -1
\end{bmatrix}
\]

\[
=
\begin{bmatrix}
13(1)+(-10)(-2) &
13(2)+(-10)(-1)
\\
20(1)+(-8)(-2) &
20(2)+(-8)(-1)
\end{bmatrix}
\]

\[
=
\begin{bmatrix}
13+20 & 26+10\\
20+16 & 40+8
\end{bmatrix}
\]

\[
=
\begin{bmatrix}
33 & 36\\
36 & 48
\end{bmatrix}
\]

\noindent
Hence,

\[
\boxed{
\operatorname{Cov}(BX^{(2)})
=
\begin{bmatrix}
33 & 36\\
36 & 48
\end{bmatrix}
}
\]

\vspace{1cm}

\noindent
\textbf{(i) \(\operatorname{Cov}(X^{(1)},X^{(2)})\)}

\vspace{0.5cm}

\noindent
\textbf{Sol.}

\vspace{0.3cm}

\noindent
The covariance matrix between \(X^{(1)}\) and \(X^{(2)}\) is the upper-right \(2\times2\) block of \(\Sigma_X\)

\vspace{0.3cm}

\noindent
Hence,

\vspace{0.3cm}

\[
\boxed{
\operatorname{Cov}(X^{(1)},X^{(2)})
=
\begin{bmatrix}
2 & 2\\
1 & 0
\end{bmatrix}
}
\]

\vspace{1cm}

\noindent
\textbf{(j) \(\operatorname{Cov}(AX^{(1)},BX^{(2)})\)}

\vspace{0.5cm}

\noindent
\textbf{Sol.}

\vspace{0.3cm}

\noindent
Using the covariance transformation property

\[
\operatorname{Cov}(AX^{(1)},BX^{(2)})
=
A\operatorname{Cov}(X^{(1)},X^{(2)})B'
\]

\[
=
\begin{bmatrix}
1 & 2
\end{bmatrix}
\begin{bmatrix}
2 & 2\\
1 & 0
\end{bmatrix}
\begin{bmatrix}
1 & 2\\
-2 & -1
\end{bmatrix}
\]

\[
A\operatorname{Cov}(X^{(1)},X^{(2)})
=
\begin{bmatrix}
1 & 2
\end{bmatrix}
\begin{bmatrix}
2 & 2\\
1 & 0
\end{bmatrix}
\]

\[
=
\begin{bmatrix}
1(2)+2(1) &
1(2)+2(0)
\end{bmatrix}
\]

\[
=
\begin{bmatrix}
4 & 2
\end{bmatrix}
\]

\noindent
Therefore,

\[
\operatorname{Cov}(AX^{(1)},BX^{(2)})
=
\begin{bmatrix}
4 & 2
\end{bmatrix}
\begin{bmatrix}
1 & 2\\
-2 & -1
\end{bmatrix}
\]

\[
=
\begin{bmatrix}
4(1)+2(-2) &
4(2)+2(-1)
\end{bmatrix}
\]

\[
=
\begin{bmatrix}
4-4 & 8-2
\end{bmatrix}
\]

\[
=
\begin{bmatrix}
0 & 6
\end{bmatrix}
\]

\noindent
Hence,

\[
\boxed{
\operatorname{Cov}(AX^{(1)},BX^{(2)})
=
\begin{bmatrix}
0 & 6
\end{bmatrix}
}
\]

\vspace{1cm}

\hrule

\vspace{0.7cm}

\noindent
\textbf{Problem 2.41 (5 points)}

\vspace{0.5cm}

\noindent
You are given the random vector

\[
X'
=
[X_1,X_2,X_3,X_4]
\]

\noindent
with mean vector

\[
\mu_X'
=
[3,2,-2,0]
\]

\noindent
and variance-covariance matrix

\[
\Sigma_X
=
\begin{bmatrix}
3 & 0 & 0 & 0\\
0 & 3 & 0 & 0\\
0 & 0 & 3 & 0\\
0 & 0 & 0 & 3
\end{bmatrix}
\]

\noindent
Let

\[
A
=
\begin{bmatrix}
1 & -1 & 0 & 0\\
1 & 1 & -2 & 0\\
1 & 1 & 1 & -3
\end{bmatrix}
\]

\vspace{0.5cm}

\noindent
\textbf{(a) Find \(E(AX)\), the mean of \(AX\)}

\vspace{0.5cm}

\noindent
\textbf{Sol.}

\vspace{0.3cm}

\noindent
Using the property

\[
E(AX)
=
AE(X)
\]

\[
E(AX)
=
\begin{bmatrix}
1 & -1 & 0 & 0\\
1 & 1 & -2 & 0\\
1 & 1 & 1 & -3
\end{bmatrix}
\begin{bmatrix}
3\\
2\\
-2\\
0
\end{bmatrix}
\]

\[
=
\begin{bmatrix}
1(3)+(-1)(2)+0(-2)+0(0)\\
1(3)+1(2)+(-2)(-2)+0(0)\\
1(3)+1(2)+1(-2)+(-3)(0)
\end{bmatrix}
\]

\[
=
\begin{bmatrix}
3-2\\
3+2+4\\
3+2-2
\end{bmatrix}
\]

\[
=
\begin{bmatrix}
1\\
9\\
3
\end{bmatrix}
\]

\noindent
Hence,

\[
\boxed{
E(AX)
=
\begin{bmatrix}
1\\
9\\
3
\end{bmatrix}
}
\]

\vspace{1cm}

\noindent
\textbf{(b) Find \(\operatorname{Cov}(AX)\), the variances and covariances of \(AX\)}

\vspace{0.5cm}

\noindent
\textbf{Sol.}

\vspace{0.3cm}

\noindent
Using the covariance transformation property

\[
\operatorname{Cov}(AX)
=
A\Sigma_XA'
\]

\noindent
Since

\[
\Sigma_X
=
3I
\]

\noindent
we have

\[
\operatorname{Cov}(AX)
=
3AA'
\]

\noindent
First calculate \(AA'\)

\[
AA'
=
\begin{bmatrix}
1 & -1 & 0 & 0\\
1 & 1 & -2 & 0\\
1 & 1 & 1 & -3
\end{bmatrix}
\begin{bmatrix}
1 & 1 & 1\\
-1 & 1 & 1\\
0 & -2 & 1\\
0 & 0 & -3
\end{bmatrix}
\]

\[
=
\begin{bmatrix}
1(1)+(-1)(-1)+0(0)+0(0)
&
1(1)+(-1)(1)+0(-2)+0(0)
&
1(1)+(-1)(1)+0(1)+0(-3)
\\
1(1)+1(-1)+(-2)(0)+0(0)
&
1(1)+1(1)+(-2)(-2)+0(0)
&
1(1)+1(1)+(-2)(1)+0(-3)
\\
1(1)+1(-1)+1(0)+(-3)(0)
&
1(1)+1(1)+1(-2)+(-3)(0)
&
1(1)+1(1)+1(1)+(-3)(-3)
\end{bmatrix}
\]

\[
=
\begin{bmatrix}
2 & 0 & 0\\
0 & 6 & 0\\
0 & 0 & 12
\end{bmatrix}
\]

\noindent
Hence,

\[
\operatorname{Cov}(AX)
=
3
\begin{bmatrix}
2 & 0 & 0\\
0 & 6 & 0\\
0 & 0 & 12
\end{bmatrix}
\]

\[
=
\begin{bmatrix}
6 & 0 & 0\\
0 & 18 & 0\\
0 & 0 & 36
\end{bmatrix}
\]

\noindent
Therefore

\[
\boxed{
\operatorname{Cov}(AX)
=
\begin{bmatrix}
6 & 0 & 0\\
0 & 18 & 0\\
0 & 0 & 36
\end{bmatrix}
}
\]

\noindent
The variances of the three linear combinations are therefore

\[
\boxed{
\operatorname{Var}(AX)_1=6
}
\]

\[
\boxed{
\operatorname{Var}(AX)_2=18
}
\]

\[
\boxed{
\operatorname{Var}(AX)_3=36
}
\]

\noindent
and all covariances between distinct linear combinations are zero

\[
\boxed{
\operatorname{Cov}(AX_i,AX_j)=0
\qquad
i\neq j
}
\]

\vspace{1cm}

\noindent
\textbf{(c) Which pairs of linear combinations have zero covariances?}

\vspace{0.5cm}

\noindent
\textbf{Sol.}

\vspace{0.3cm}

\noindent
Let the three linear combinations be

\[
Y_1=X_1-X_2
\]

\[
Y_2=X_1+X_2-2X_3
\]

\[
Y_3=X_1+X_2+X_3-3X_4
\]

\noindent
From the covariance matrix obtained in part (b)

\[
\operatorname{Cov}(AX)
=
\begin{bmatrix}
6 & 0 & 0\\
0 & 18 & 0\\
0 & 0 & 36
\end{bmatrix}
\]

\noindent
we have

\[
\operatorname{Cov}(Y_1,Y_2)=0
\]

\[
\operatorname{Cov}(Y_1,Y_3)=0
\]

\[
\operatorname{Cov}(Y_2,Y_3)=0
\]

\noindent
Therefore all pairs of linear combinations have zero covariance

\[
\boxed{
(Y_1,Y_2)
\qquad
(Y_1,Y_3)
\qquad
(Y_2,Y_3)
}
\]

\end{document}